\section{The probabilistic interpretation and the Li-Yau inequality}

The main goal of this section is to prove the following two results 

\begin{proposition}[Li-Yau-Wang application] \label{prop:ultracon_app}
Let $V$ be a convex potential in $\R$ satisfying Conditions \ref{cond:mass} and \ref{cond:coarse_flow}. Let $(P_t)_{t \ge 0}$ be the semigroup associated to $-\Delta + \nabla V \cdot \nabla$. Then for all $d\ge 1$ one has the estimate
$$
\|P_{s+t}\|_{L^2\to L^\infty} \le C_{LYW}(s,t)
$$
with $C_{LYW}(s,t) = (2 \pi s)^{-1/2}  \exp(\frac{5}{s})\frac{s}{s-t}$
\end{proposition}

\begin{proposition}[Good approximation by flow] \label{prop:approx_by_flow}
    If \( f \) is \( C \)-Lipschitz, then for all \( t \ge 0 \) it holds that $P_t f (x) \le f(\Phi_t(x)) + C 2\sqrt{t}$.
\end{proposition}

Our approach is rooted in the probabilistic representation of the semigroup, which we explore first. We will construct the semigroup from a specific stochastic differential equation and then derive its fundamental properties, such as its invariant measure and contractivity. These properties will then be used to prove the two propositions.


\subsection{Probabilistic representation of parabolic semigroups} \label{sec:probabilistic}

We present a summary of the tools utilized in the probabilistic theory of semigroups, drawing on \cite{bakry_analysis_2014} and \cite{chen_lecture_2025} as references. We begin by defining the semigroup associated with a general Markov process.

\begin{definition}
    Let $(X_t)_{t\ge 0}$ be a Markov process. The associcated \emph{semigroup} acts on a suitable function $f$ by
    $$
    P_t f(x) := \E[f(X_t) | X_0 = x ].
    $$
    The \emph{infinitesimal generator} $L$ of this semigroup is defined by the limit
    $$ 
    Lf := \lim_{t \to 0^+} \frac{P_t f - f}{t}, 
    $$
    for all functions $f$ for which this limit exists.
\end{definition}

The generator is fundamental because it connects the semigroup to a partial differential equation, governing its evolution in time.

\begin{proposition} \label{prop:partial_P}
    The generator $L$ of a semigroup $(P_t)_{t \ge 0}$ satisfies the relation $\partial_t P_t = LP_t = P_t L$.
\end{proposition}
\begin{proof}
    The semigroup property, $P_{t+s} = P_t P_s$, allows us to write
    $$
    L P_t f 
    = \lim_{s \to 0^+} \frac{P_s P_t f - P_t f}{s}
    = \lim_{s \to 0^+} \frac{P_{t+s} - P_t}{s} f = \partial_t P_t f.
    $$
    By interchanging $P_t$ and the limit, which is justified by the continuity of $P_t$, we also obtain
    $$
    P_t L f 
    = P_t \lim_{s \to 0^+} \frac{P_s f - f}{s}
    = \lim_{s \to 0^+} \frac{P_{t+s} f - P_t f}{s}
    $$
\end{proof}

The specific semigroup of interest for our problem is the semigroup generated by the operator $L_V = -\Delta + \nabla V \cdot \nabla$.  The Markov process associated to this semigroup is the solution of the following stochastic differential equation (SDE):

\begin{definition} \label{def:langevin_process}
    We define the Langevin diffusion $(X_t)_{t \ge 0}$ as the Markov process solving the SDE
    $$
    \dd{X_t} = - \nabla V(X_t) \dd{t} + \sqrt{2} \dd{B_t},
    $$
    where $(B_t)_{t \ge 0}$ is a Brownian motion in $\R^2$
\end{definition}

% Remark: Chen has a source for the existence of a unique solution if V is lipschitz or something like that

\begin{proposition}
    $L_V = - \Delta + \nabla V \cdot \nabla$ is the generator of the semigroup $(P_t)_{t \ge 0}$ associated to the Markov process which solves the stochastic differential equation $\dd{X_t} = - \nabla V(X_t) \dd{t} + \sqrt{2} \dd{B_t}$
\end{proposition}
\begin{proof}
    By Itô's lemma, for a smooth function $f$, the differential $\dd{f}(X_t)$ is given by
    \begin{align*}
    \dd{f(X_t)} &= \left( \pdv{f}{t} - \nabla V \cdot \nabla f + \frac{1}{2} \operatorname{Tr}(\sqrt{2} D^2 f \sqrt{2}) \right)\dd{t} + \nabla f \cdot \sqrt{2}\dd{B_t} \\
    &= -\left(\nabla V \cdot \nabla -  \Delta \right) f \dd{t} + \nabla f \cdot \sqrt{2}\dd{B_t}
    \end{align*}
    Taking the expectation conditioned on $X_0=x$ and noting that the Brownian motion term has mean zero, we get
    \[
    \E[\dd{f}(X_t)\mid X_0=x]
    = \E\big[(\Delta - \nabla V \cdot \nabla)f(X_t)\mid X_0=x\big]\dd{t}.
    \]
    Hence, for the semigroup $P_t f(x)=\E[f(X_t)\mid X_0=x]$,
    \[
    \partial_t P_t f(x) = P_t\big(\Delta - \nabla V \cdot \nabla\big)f(x).
    \]
    Therefore, the infinitestimal generator is $L=\Delta-\nabla V\cdot\nabla$.
\end{proof}

\begin{corollary} \label{prop:flow_of_eig}
    If $\phi$ is an eigenfunction of the generator $L$ with eigenvalue $\lambda$, then $\partial_t P_t \phi = \lambda P_t \phi$. This implies that $\phi = e^{\lambda t}P_t \phi$.
\end{corollary}

An important property of the semigroup is its invariant measure.

\begin{definition}[Invariant measure]
    We say that a measure $\mu$ is invariant for $P$ if $\forall t \ge 0 : \int P_t f \dd{\mu} = \int f \dd{\mu}$
\end{definition}

\begin{theorem}[Properties of the Semigroup] \label{thm:prop_of_semigroup}
The Langevin semigroup $(P_t)_{t\ge 0}$ has the following properties:
\begin{enumerate}
    \item \textbf{Maximum Principle:} If $f \le g$ pointwise, then $P_t f \le P_t g$ pointwise.
    \item \textbf{$L^p$-Contractivity:} If $\mu$ is invariant for $(P_t)_{t \ge 0}$, then for any $p \in [1, \infty]$, $P_t$ is a contraction on $L^p(\mu)$, i.e., $\|P_t f\|_{L^p(\mu)} \le \|f\|_{L^p(\mu)}$.
    \item \textbf{Gradient Estimate:} For a differentiable function $f$, $|\nabla P_t f| \le P_t(|\nabla f|)$ pointwise.
    \item \textbf{Lipschitz Preservation:} If $f$ is $C$-Lipschitz, then $P_t f$ is also $C$-Lipschitz for all $t \ge 0$.
\end{enumerate}
\end{theorem}
\begin{proof} \,
\begin{enumerate}
\item \textbf{Maximum Principle:} This is a direct consequence of the definition and the monotonicity of expectation.
    \[ P_t f(x) = \E[f(X_t) | X_0 = x] \le \E[g(X_t) | X_0 = x] = P_t g(x). \]

\item \textbf{$L^p$-Contractivity:} For $p=\infty$, the result is by definition of $P_t$ as an expected value. For $1 \le p < \infty$, Jensen's inequality gives the pointwise estimate $|P_t f|^p = |\E[f(X_t)]|^p \le \E[|f(X_t)|^p] = P_t(|f|^p)$. Integrating with respect to the invariant measure $\mu$ yields
    \[
        \|P_t f\|_{L^p(\mu)}^p
        = \int |P_t f|^p \dd{\mu}
        \le \int P_t(|f|^p) \dd{\mu}
        = \int |f|^p \dd{\mu}
        = \|f\|_{L^p(\mu)}^p.
    \]

 \item \textbf{Gradient Estimate:} This follows from the fact that $P_t$ is an expected value; $P_t$ commutes with $\nabla$ therefore $|\nabla P_t f| = |P_t(\nabla f)|$, followed by an application Jensen's inequality to the absolute value.

\item \textbf{Lipschitz Preservation:} Let $f$ be $C$-Lipschitz. We can approximate $f$ by a sequence of smooth, $C$-Lipschitz functions $f_n$. For each $f_n$, the gradient estimate gives $\|\nabla (P_t f_n)\|_\infty \le \|P_t(|\nabla f_n|)\|_\infty \le \|P_t C\|_\infty = C$. Thus, $P_t f_n$ is also $C$-Lipschitz. Since $P_t$ is a continuous operator on $L^p$, $P_t f_n \to P_t f$ pointwise. Then \[ |P_t f(x) - P_t f(y)| = \lim_{n\to\infty} |P_t f_n(x) - P_t f_n(y)| \le C|x-y|, \] for any $x,y$, showing that $P_t f$ is $C$-Lipschitz.
\end{enumerate}
\end{proof}

For our purposes we need the above properties for the Langevin semigroup, which follow once we identify the invariant measure.

\begin{lemma} \label{lem:measure_is_invariant}
    The measure \( \mu_d \sim e^{-V} \) is invariant for the Langevin semigroup \( (P_t)_{t \ge 0} \).
\end{lemma}
\begin{proof}
    $$
    \int ((-\Delta + \nabla V \cdot \nabla) f) \, e^{-V} \dd{x}
    = \int (\nabla V \cdot \nabla f) \, e^{-V} \dd{x} - \int (\Delta f) \, e^{-V} \dd{x}
    $$
    integrating by parts the second term
    $$
    \int ((-\Delta + \nabla V \cdot \nabla) f) \, e^{-V} \dd{x} = \int (\nabla V \cdot \nabla f) \, e^{-V} \dd{x} + \int (\nabla V \cdot \nabla f) \, e^{-V} \dd{x} = 0.
    $$
    In, particular
    \[
        \dv{t} \int P_t f \dd{\mu}
        = \int (-\Delta + \nabla V \cdot \nabla) P_t f \dd{\mu}
        = 0,
    \]
    hence $\int P_t f \dd{\mu}$ does not depend on $t$, in particular $\int P_t f \dd{\mu} = \int P_0 f \dd{\mu} = \int f \dd{\mu}$.
\end{proof}

\begin{lemma}[Corollary 1 in \cite{chen_lecture_2025}] \label{lem:Bakry-Emery-Poincare}
    Let \( \rho > 0 \), and let \( \mu = e^{-V} \) be a probability measure. Assume \( \nabla^2 V \succeq \rho \mathbb{I}_n \). Then \( \mu \) satisfies the Poincaré inequality
    \begin{equation*}
        \operatorname{Var}_\mu(f) \le 1/\rho \int_{\mathbb{R}^n} |\nabla f|^2 \dd\mu.
    \end{equation*}
    If we assume that \( f \) has mean zero, then \( \operatorname{Var}_\mu(f) = \int_{\mathbb{R}^n} |f|^2 \dd\mu \) and therefore \( \lambda^{(1)}_\mu \ge \rho \).
\end{lemma}

\subsection{Proof of good approximation by flow}

\begin{proof}[Proof of \Cref{prop:approx_by_flow}]
We claim that $\E[\|X_t - \Phi_t\|^2] \le 4t$. Following this, the argument proceeds as:  
$P_t f(X_0) = \E[f(X_t)] = \E[f(\Phi_t)] + \E[ f(X_t) - f(\Phi_t) ]$. Since $\Phi_t$ is deterministic, we can drop the expected value. Given that $f$ is $C$-Lipschitz, we arrive at $P_t f(X_0) \le f(\Phi_t) + C \E[|X_t - \Phi_t|]$. Applying Jensen's inequality, it follows that $\E[\|X_t - \Phi_t\|] \le \sqrt{\E[\|X_t - \Phi_t\|^2]}$. Incorporating our claim, we derive
$$
P_t f(X_0) \le  f(\Phi_t) + C 2\sqrt{t}.
$$

    To prove the claim let \( Y_t := X_t - \Phi_t \). The differential of \( Y_t \) is given by \( \dd{Y_t} = -(\nabla V(X_t) - \nabla V(\Phi_t)) \dd{t} + \sqrt{2}\dd{B_t} \).
    By Itô's lemma, we have
    \begin{equation*}
        \dd{\|Y_t\|^2} = -2(\nabla V(X_t) - \nabla V(\Phi_t)) \cdot Y_t \dd{t} + 4 \dd{t} + 2 Y_t \cdot \sqrt{2}\dd{B_t}.
    \end{equation*}
    In expectation, we obtain
    \begin{equation*}
        \dd{\E[\|Y_t\|^2]} = - 2(\nabla V(X_t) - \nabla V(\Phi_t)) \cdot Y_t \dd{t} + 4\dd{t}.
    \end{equation*}
    Since \( V \) is convex, we have that \( \nabla V(X_t) - \nabla V(\Phi_t) \cdot (X_t-\Phi_t) \ge 0 \) and therefore $\dd{\E[\|Y_t\|^2]} \le 4 \dd{t}$, from which it immediately follows that \( \E[\|X_t - \Phi_t\|^2] \le 4t \).
\end{proof}

\subsection{Ultracontractivity estimates}

Most classical approaches to PDE regularity theory (such as those based on de Giorgi or Moser iterations) are based on comparing the solutions to the PDE at two adjacent scales (for example comparing an $L^p$ norm of the solution on a certain ball to an $L^q$ norm on a larger ball, for $q<p$). These estimates have poor dimensional dependency, because they absorb volume doubling factors for the ball (of the form $2^d$). Similarly, classical estimates on a fixed dimension for elliptic equations of the form $-\Delta + w(x) \cdot \nabla $ treat $w$ as a perturbation, and the constants on the bounds degenerate as $w$ grows to infinity. Both of these facts are not artifacts of the proof methods for general operators.

In our case, however, we are taking $w = (d-1)\nabla V_0$. This is one of the few notable exceptions where one can keep the constants in the regularity estimates for PDE uniformly bounded, both in the dimension, and the regularity of the coefficients of the first order term. This fact is known as Wang's Harnack inequality.

\begin{theorem}[Wang's Harnack inequality, Lemma 2.1 in \cite{wang_logarithmic_1997}] \label{thm:wangs_harnack}
    Let  \( f \in C_b(S) \). For any \( t \ge 0\) and positive \( g \in C[0,t] \), we have
    \[
        |P_t f(x)|^2 \le \exp\left(\frac{|x-y|^2}{2t}\right) P_t |f|^2(y),
    \]
    for all \( x,y \in S\).
\end{theorem}

As a corollary we have

\begin{theorem}[Ultracontractivity \cite{rockner_supercontractivity_2003}] \label{thm:ultracon}
    If \( (\int \exp\left(-\frac{|x|^2}{2s}\right) \dd{\mu})^{-1} \le C_1 \) and \( \|P_{t} \exp\left(\frac{|x|^2}{4s}\right)\|_{L^{\infty}} \le C_2 \) , then
    \[
        \|P_{s+t} f\|_{L^{\infty}} \le C_1 C_2 \|f\|_{L^2}.
    \]
\end{theorem}
\begin{proof}
    It suffices to prove that
    \begin{equation}
        |P_{s} f|(x) \le C_1 \| f \|_{L^2} \exp\left(\frac{|x|^2}{4s}\right) .
    \end{equation}
    Indeed, using the maximum principle, we have
    \begin{equation*}
        \|P_{s+t} f\|_{L^{\infty}} 
        = \|P_{t} P_{s} f\|_{L^{\infty}} 
        \le C_1 \|f\|_{L^2} \|P_{t} \exp\left(\frac{|x|^2}{4s}\right)\|_{L^{\infty}} 
        \le C_1 \|f\|_{L^2} C_2,
    \end{equation*}
    and we are done.
    According to \Cref{lem:measure_is_invariant}, \( \int P_{s} |f|^2(y) \dd{\mu(y)} = \int |f|^2(y) \dd{\mu(y)} = \| f \|_{L^2}^2 \).
    By integrating the Wang's Harnack inequality (\Cref{thm:wangs_harnack}), we have
    \begin{equation}
        |P_{s} f|^2(x) \int \exp\left(-\frac{|x-y|^2}{2s}\right) \dd{\mu(y)}
        \le \int P_{s} |f|^2(y) \dd{\mu(y)}
        = \| f \|_{L^2}^2,
    \end{equation}
    which can be rearranged to
    \begin{align*}
        |P_{s} f|^2(x) 
        &\le \left( \int \exp\left(-\frac{|x-y|^2}{2s}\right) \dd{\mu(y)} \right)^{-1} \| f \|_{L^2}^2 \\
        &\le \underbrace{\left( \int \exp\left(-\frac{|y|^2}{2s}\right) \dd{\mu(y)} \right)^{-1}}_{\le C_1} \exp\left(\frac{|x|^2}{2s}\right) \| f \|_{L^2}^2
    \end{align*}
\end{proof}

\subsubsection*{Estimates for the Ultracontractivity constants}

This section is dedicated to giving estimates of the constants in \Cref{thm:ultracon} for the Lebesgue measure weighted by $e^{-V}$. From these two constant, \Cref{prop:ultracon_app} will immediately follow.

\begin{lemma}
    Let $s > 0$. Assuming that $V$ satisfies \Cref{cond:mass}, we have 
    \[
        \left(\int \exp\left(-\frac{|x|^2}{2s}\right) \dd{\mu}\right)^{-1} \le C_1 \approx e^{2/s} (2\pi s)^{-1/2}.
    \]
\end{lemma}
\begin{proof}
%    \begin{align*}
%        \frac{1}{Z_d} \int \int e^{-(d-1)\frac{x^2}{2\tau_1}} e^{-(d-1)(\frac{y^2}{2\tau_1}+\frac{y}{2}-\log y)} \dd{x}\dd{y}
%    \end{align*}
Under \Cref{cond:mass} the measure $\mu_d$ converges to the uniform measure on \( [-\pi/2,\pi/2] \times \{2\} \), therefore, as $d \to +\infty$, 
    \begin{equation}
        \int \exp\left(-\frac{|x|^2}{2s}\right) \dd{\mu} \to \int_{-1}^1 e^{-\frac{x^2+4}{2s}} \dd{x} = e^{-2/s}\sqrt{2\pi s} \erf{\frac{\pi}{2 \sqrt{2s}}}.
    \end{equation}
    \begin{equation}
        C_1 \approx e^{2/s} (2\pi s)^{-1/2}
    \end{equation}
\end{proof}


\begin{lemma}\label{lem:x_sq_le}
    For any \( x,y \in \R^2 \) we have $|x|^2 \le 2|x-y|^2 + 2|y|^2$
\end{lemma}
\begin{proof}
    Write \( x = (x-y) + y \). By the triangle inequality,
    \[
        |x|^2 \le (|x-y| + |y|)^2 = |x-y|^2 + 2|x-y| |y| + |y|^2.
    \]
    Using that \( 2|x-y||y| \le |x-y|^2 + |y|^2 \), we obtain
    \[
        |x-y|^2 + 2|x-y| |y| + |y|^2 \le 2|x-y|^2 + 2|y|^2,
    \]
    which proves the claim.
\end{proof}

\begin{lemma} \label{lem:coupling}
    Let \( \mu_t^{x_0} \) be the law of the Langevin process $(X_t)_{t \ge 0}$ (\Cref{def:langevin_process}) starting at $X_0 = x_0$. There is a coupling measure $\pi$ of \( \mu_t^{x_0} \) and \( \gamma \sim \mathcal N(0,tI_2) \) such that almost surely
    \[
        |X_t - \Phi_t(X_0)| \le |\gamma|.
    \]
\end{lemma}
\begin{proof}
    Let $Y_t:= X_t-\Phi_t(X_0)$, and let $V_t(x):= V(x-\Phi_t(x_0))-\nabla V(\Phi_{t}(x_0)\cdot x$. Then $V_t$ is a convex potential with minimum at $x=0$, and $Y_t$ satisfies the SDE
    $$
    \begin{cases}
        d Y_t = d B_t^{(2)} - \nabla V_t(Y_t) dt\\
        Y_0 = 0
    \end{cases}
    $$
    where we are using the super-index in $d B_t^{(2)}$ to denote that it is a two-dimensional Brownian motion. 
    
    The function $Y\mapsto |Y|$ in $\mathbb R^2$ has gradient equal to  $Y/|Y|$ and Laplacian equal to $\frac{1}{|Y|}$. In particular $|Y_t|^2$ satisfies the SDE
    $$
    d |Y_t| =  d \tilde B_t +\frac{1}{2|Y_t|} dt - \langle Y , \nabla V_t(Y_t)\rangle dt.
    $$
    
    The process, $B^{(2)}_t$ has, at time $t$ the same law as a Gaussian of variance $t I_2$, and $|B^{(2)}_t|^2$ satisfies the SDE
    $$
    d |B^{(2)}_t| = d \tilde {\tilde B}_t +\frac{1}{2|B^{(2)}_t|} dt.
    $$
    We can consider two versions of $|Y_t|^2$ and $|B^{(2)}_t|^2$ where we couple the driving Brownian motions $\tilde {\tilde B}_t$ and ${\tilde B}_t$. For this coupling, another application of Ito's formula leads to
    $$
    d (|Y_t| -|B^{(2)}_t|)  = \frac{|Y_t|-|B_t|}{|Y_t|+|B_t^{(2)}|} - \langle Y_t/|Y_t|,\nabla V_t(Y_t)\rangle.
    $$
    The denominator $|Y_t|+|B_t^{(2)}|$ is almost surely bounded far from zero for finite times by the non-recurrence of two dimensional Brownian motion, which, by an application of Gronwall's inequality, shows that $|Y_t| -|B^{(2)}_t| \le 0$. In particular, there is a coupling of $Y_t$ and $B_t^2$ for which $|Y_t| \le |B^{(2)}_t|$ almost surely. Unfolding the definitions shows the result.
    
\end{proof}


\begin{lemma}
    Let $s,t > 0$, assuming that $V$ satisfies \Cref{cond:coarse_flow}, we have
    \[ 
        \left\|P_t \exp\left(\frac{|x|^2}{4s}\right)\right\|_{L^{\infty}} \le e^{5/s} \frac{s}{s-t}.
    \]
\end{lemma}
\begin{proof}
    Let \( x_0 \in S \).
    Applying Lemma~\ref{lem:x_sq_le} to \( y = \Phi_{t}(x_0) \) we obtain
    \begin{align*}
        \left[P_{t} \exp(\frac{|x|^2}{4s})\right](x_0) 
        &= \E_{x \sim \mu_{t}^{x_0}} \left[ \exp(\frac{|x|^2}{4s}) \right] \\
        &\le \E_{x \sim \mu_{t}^{x_0}} \left[ \exp(\frac{2|x-\Phi_{t}(x_0)|^2 + 2|\Phi_{t}(x_0)|^2}{4s}) \right] \\
        &= \exp(\frac{|\Phi_{t}(x_0)|^2}{2s}) \E_{x \sim \mu_{t}^{x_0}} \left[ \exp(\frac{|x-\Phi_{t}(x_0)|^2}{2 s}) \right] \\
        &= \exp(\frac{|\Phi_{t}(x_0)|^2}{2s}) \E_{(x,\gamma) \sim \pi}\left[\exp(\frac{|x-\Phi_{t}(x_0)|^2}{2 s})\right] \\
        &\le \exp(\frac{|\Phi_{t}(x_0)|^2}{2s}) \E_{(x,\gamma) \sim \pi}\left[\exp(\frac{|\gamma|^2}{2 s})\right]
    \end{align*}
    The expected value is a Gaussian integral, therefore
    \begin{align*}
        E_{(x,\gamma) \sim Coupling}[\exp(\frac{|\gamma|^2}{2 s}) ]
        &= \frac{1}{2\pi t} \int_{\R^2} \exp(\frac{|z|^2}{2 s} -\frac{|z|^2}{2t}) \dd{z} \\
        &= \frac{1}{2\pi t}\int_{\R^2} \exp(-|z|^2(\frac{1}{2 t} -\frac{1}{2s})) \dd{z} \\
        &= \frac{s}{s - t}.
    \end{align*}
\end{proof}

\subsubsection*{Optimal times for ultracontractivity}

For our purposes we need to apply \Cref{prop:ultracon_app} to functions  of the form $e^{\lambda (s+t/d)}P_{s+t/d}$. The question that arises is how do we choose $s$ and $t$ such that the resulting constant $C_{WLY}$ is small. Recall that $C_{LYW}(s,t/d) = (2 \pi s)^{-1/2}  \exp(\frac{5}{s})\frac{s}{s-t/d} \approx (2 \pi s)^{-1/2}  \exp(\frac{5}{s})$ for large enough $d$, hence we should optimize $e^{\lambda s +5/s} s^{-1/2}$. The logarithm of this expression is $\lambda s + 5/s - \frac{1}{2}\log s$, which is minimal at $s = \frac{1 + \sqrt{1+80\lambda}}{4 \lambda}$. Evaluated at the first eigenvalue $\lambda=1$ this equals $s=2.5$.